\chapter{Língua Portuguesa}
\label{cap:portugues}

\begin{edital}
Compreensão e interpretação de textos; tipos e gêneros textuais; ortografia;
coesão (referenciação, conectores, tempos/modos verbais); estrutura
morfossintática (classes de palavras, coordenação, subordinação, pontuação,
concordância verbal e nominal, regência, crase, colocação pronominal);
reescrita e significação.
\end{edital}
\peso{12 questões, peso 1 --- mesmo peso do Inglês. Interpretação é o eixo. Na
Dataprev 2024 foram as questões Q1--Q12; o carro-chefe é gramática aplicada e
interpretação de frase curta, não redação nem texto longo.}

\section{Onde a prova se concentra}

A FGV mistura \textbf{interpretação} com \textbf{gramática aplicada} (não
gramática decorada e solta). Prioridades:

\begin{enumerate}
\item \textbf{Interpretação e reescrita} (o maior naipe): sentido do texto,
do trecho, de conectivos; substituição de palavras/trechos mantendo o
sentido.
\item \textbf{Sintaxe do período:} coordenação $\times$ subordinação, função
dos termos (sujeito, adjunto adnominal $\times$ complemento nominal), orações
reduzidas.
\item \textbf{Concordância e regência:} verbal e nominal; crase.
\item \textbf{Coesão:} referenciação (pronomes, elipse), conectores e o que
eles expressam (causa, consequência, concessão, condição).
\item \textbf{Pontuação:} justificar o uso da vírgula (aposto, adjunto
deslocado, oração subordinada, enumeração).
\end{enumerate}

\section{Pares e conceitos que a FGV cobra}

\begin{itemize}
\item \textbf{Adjunto adnominal} (liga-se a substantivo, indica
posse/qualidade) $\times$ \textbf{complemento nominal} (completa o sentido
de nome, é alvo da ação).
\item \textbf{Oração subordinada substantiva} (função de substantivo:
sujeito, objeto) $\times$ \textbf{adjetiva} (função de adjetivo, restritiva
$\times$ explicativa) $\times$ \textbf{adverbial}.
\item \textbf{Oração reduzida} (verbo no infinitivo/gerúndio/particípio, sem
conectivo) --- desenvolvê-la revela a função.
\item \textbf{Regência:} ``obedecer \textbf{a}'', ``assistir \textbf{a}''
(ver), ``lembrar-se \textbf{de}'', ``residir \textbf{na}/à'' --- verbos que a
FGV adora com regência trocada.
\item \textbf{Concordância:} ``anexo/anexa'' concorda; ``é
proibido/proibida a entrada'' (varia com o determinante); ``meio'' advérbio é
invariável (``meio confusa'').
\item \textbf{Discurso direto $\times$ indireto:} o direto reproduz a fala
(aspas, travessão, verbo de elocução); não ``demarca a soberania do
narrador''.
\end{itemize}

\begin{pegadinha}
O enunciado é quase sempre no \textbf{NEGATIVO}: ``assinale a opção
INCORRETA'', ``em que o termo NÃO funciona como...''. Metade dos erros é ler
correto onde está incorreto --- circule a palavra-chave antes de olhar as
alternativas.

Em interpretação, o distrator é a leitura que \textbf{EXTRAPOLA} ou
\textbf{INVERTE} a tese. Ex.: num verso de Pessoa, ``visão negativa do fazer
poético'' ou ``torna o poeta imune à dor'' --- o texto não diz nada disso; o
correto é o mais literal (``o fingimento é matéria de poesia'').

Distrator com termo \textbf{absoluto}: ``É impossível ser feliz e viver'',
``Viver pressupõe felicidade''. A frase original só descrevia um paradoxo ---
desconfie de qualquer opção que radicaliza.

Em gramática, a alternativa errada costuma ter \textbf{um erro pontual
plantado} (regência trocada, concordância de ``meio/anexo/proibido'', crase
indevida). O resto da frase parece natural de propósito.
\end{pegadinha}

\begin{comosair}
Volte ao texto e teste cada alternativa contra o que está \textbf{ESCRITO},
não contra o que ``faz sentido no mundo''. Se a opção precisa de informação
que o trecho não dá, ela extrapola: elimine. Elimine primeiro os absolutos
(sempre, nunca, impossível, todos) e as leituras que agregam juízo de valor
ausente. Revise as regras-armadilha clássicas: concordância de
``anexo/proibido/meio'' (variam ou não conforme o caso), regência de
obedecer/simpatizar/residir/lembrar, e o uso da vírgula (aposto, adjunto
deslocado, oração intercalada). Em oração subordinada, saiba nomear
rapidamente: substantiva (função de substantivo --- sujeito, objeto),
adjetiva (delimita/explica um nome), reduzida (verbo no
infinitivo/gerúndio/particípio, sem conectivo).
\end{comosair}

\begin{jacaiu}
Análise sintática de frase curta (``É preciso estar atento e forte'');
adjunto adnominal; regência verbal; concordância nominal; interpretação de
verso (Pessoa, Gonzaga); justificativa da vírgula; discurso direto. Rode
\texttt{./quiz.py portugues}.
\end{jacaiu}

\section{Pontos de atenção / pesquisa extra}

\begin{itemize}
\item \textbf{Crase:} casos obrigatórios (locuções femininas, ``à medida
que''), proibidos (antes de verbo, masculino, pronomes) e facultativos.
\item \textbf{Colocação pronominal:} próclise (palavra atrativa: negação,
pronome relativo, advérbio), ênclise, mesóclise.
\end{itemize}
